\newcommand{\hdeltaset}{\widehat{\deltaset}}
\newcommand{\hbfG}{\hat{\bfG}}
\newcommand{\kk}{{k \times k}}
\section{Analysis for \cref{sec:sensitivity}}
\label{sec:startofproofs}

\newcommand{\Gsq}{\bfG_{\sf sq}}
\renewcommand{\g}[2]{g\idx{#1}{#2}}
\newcommand{\cond}{\emph{(1)-(3)}\,}
\subsection{From scalar to vector contributions}
\begin{thm}\label{thm:multidim}
Let $\bfC \in \R^{\dimension \times k}$ which satisfies 
\[
\| \bfC \contrib \|_2 \le 1 \qquad \forall \contrib \in \{ -1, 1\}^k,
\]
and let $\bfG \in \R^{k \times \mdim}$ such that each row $\bfG\idx{i}{:}$ for $i \in [k]$ satisfies $\|\bfG\idx{i}{:}\|_2 \le 1$. Suppose at least one of the following conditions hold:
\begin{enumerate}
    \item We have $k = 1$ or $k = 2$ participations.
    \item All the entries of $\bfC^\top\bfC$ are non-negative.
    \item The rows of $\bfG$ are all co-linear, $\bfG_{[i,:]} = u_i\cdot  \bfG_{[1,:]}$ for $u_i \in \{-1, 1\}$, $i > 1$.
    \item The rows of $\bfG$ are all orthogonal, $\ip{\bfG_{[i,:]}}{\bfG_{[j,:]}} = 0$, $\forall i \neq j$, $i, j \in [k]$.
\end{enumerate}
Then,
\begin{equation}
\lfrob{ \bfC \bfG } \le 1.
\label{eq:th143}
\end{equation}
Furthermore, the following statements are also true without assuming conditions~\cond above.
\begin{itemize}
    \item $\lfrob{ \bfC \bfG } \leq \frac{\pi}{2}$.
    \item If we replace the condition on $\bfG$ to $\forall i \in [k], \|\bfG\idx{i}{:}\|_1 \le 1$, then $\lfrob{\bfC \bfG } \le 1$.
\end{itemize}
\label{thm:901r}
\end{thm}
\mypar{Note} \cref{thm:multidim} is generally applied to $\bfC\idx{:}{\pi}$, the sub-matrix of some $\bfC \in \R^\dimdim$ formed by keeping only columns selected by a particular participation pattern $\pi$.

\begin{proof}[Proof \cref{thm:901r}]
Let $\bfC=\begin{bmatrix} \bfc_1 & \bfc_2 & & \cdots & \bfc_k \end{bmatrix} $ with each $\bfc_i\in \mathbb{R}^n$ being a column vector. Also let we write $\g{i}{j}$ for the $(i,j)$-th entry of $\bfG$. It will also be useful to note
\begin{equation}\label{eq:traceCG}
\lfrob{\bfC\bfG}^2 
  = \tr\big(\bfC \bfG (\bfC \bfG)\tp \big)
  = \tr\big(\bfC\tp \bfC \bfG\bfG\tp \big).
\end{equation}
In the following we prove each of the individual cases of Theorem~\ref{thm:901r}.

\mypar{When $k=1$ or $k=2$:} For $k=1$, we have
\begin{equation*}
    \lfrob{\bfC\bfG}^2=\ltwo{\bfc_1}^2\left(\sum\limits_{j=1}^d\g{1}{j}\right)^2\leq\ltwo{\bfc_1}^2=\max\limits_{u\in\{\pm 1\}}\ltwo{u\cdot\bfc_1}^2.%
\end{equation*}
An equivalent argument is used in \citet[Thm. 3.1]{denisov22matfact}.

For $k=2$, we have the following:
\begin{align}
    \lfrob{\bfC\bfG}^2&=\sum_{i=1}^2\ltwo{\bfc_i}^2\cdot\left(\sum_{j=1}^d\g{i}{j}^2\right) 
         + \left(2\g{1}{1}\g{2}{1}\ip{\bfc_1}{\bfc_2}\right)
         + \cdots +
         {\left(2\g{1}{d}\g{2}{d}\ip{\bfc_1}{\bfc_2}\right)}\nonumber\\
         &\leq \sum_{i=1}^2\ltwo{\bfc_i}^2+\left(2|\g{1}{1}|\cdot|\g{2}{1}||\ip{\bfc_1}{\bfc_2}|\right)+\cdots+\left(2|\g{1}{d}|\cdot|\g{2}{d}||\ip{\bfc_1}{\bfc_2}|\right)\nonumber\\
         &\leq \sum_{i=1}^2\ltwo{\bfc_i}^2+\left(\g{1}{1}^2+\g{2}{1}^2\right)|\ip{\bfc_1}{\bfc_2}|+\cdots+\left(\g{1}{d}^2+\g{2}{d}^2\right)|\ip{\bfc_1}{\bfc_2}| \label{eq:p9013} \\
         &=\ltwo{\bfc_1}^2+\ltwo{\bfc_2}^2+2|\ip{\bfc_1}{\bfc_2}|
          =\max\left\{\left(\bfc_1+\bfc_2\right)^2,\left(\bfc_1-\bfc_2\right)^2\right\} \nonumber\\
          &\leq \max\limits_{\bfu\in\{\pm 1\}^2}\ltwo{\bfC\bfu}^2\leq 1,\nonumber %
\end{align}
where \cref{eq:p9013} follows from the standard $\texttt{A.M.}\geq\texttt{G.M.}$ inequality. %

\mypar{All the entries of $\bfC^\top\bfC$ are non-negative:} Let $\bfX = \bfC\tp\bfC$ and $\hbfG = \bfG \bfG\tp$. Observe  $\hbfG\idx{i}{j}\in [-1, 1]$, and using \cref{eq:traceCG}, when $\bfX$ is elementwise non-negative, $\tr\big(\bfX \hbfG\big)$ is maximized when $\hbfG = \mathbf{1}^{k \times k} = \widehat{\bfu} \widehat{\bfu}\tp$ by choosing $\widehat{\bfu} = \mathbf{1}^k$. Hence,
\begin{equation}
    \lfrob{\bfC\bfG}\leq \tr\big(\bfC^\top\bfC\widehat{\bfu} \widehat{\bfu}\tp\big)=\ltwo{\bfC\widehat{\bfu}}\leq \max\limits_{\bfu\in\{\pm 1\}^k}\ltwo{\bfC\bfu}^2.\label{eq:p9001}
\end{equation}
\cref{eq:p9001} completes the proof.

\mypar{The rows of $\bfG$ are co-linear:}
By the convexity of $\lfrob{\bfC\bfG}^2$ with respect to the matrix $\bfG$, we  may assume the rows of $\bfG$ are of $\ell_2$ norm 1. Under the colinearity assumption, this translates to $\bfG_{[i,:]} = u_i \bfG_{[1,:]}$, with each $u_i\in\{\pm 1\}$. Let $\bfu = [u_1, \dots, u_k] \in \{-1, 1\}^k$. Then for the matrix $\bfG\bfG^\top$ we have the following:
\begin{equation*}    
  [\bfG \bfG\tp]\idx{i}{j} 
    = \ip{\bfG_{[i,:]}}{\bfG_{[j,:]}} 
    = u_i u_j \ip{\bfG_{[1,:]}}{\bfG_{[1,:]}} = u_i u_j.%
\end{equation*}
which implies
\begin{equation}\label{eq:GGuu}
    \bfG \bfG\tp = \bfu \bfu\tp.
\end{equation}
Using \cref{eq:GGuu} with \cref{eq:traceCG}, we have
\[
\lfrob{\bfC\bfG}^2 
  = \tr\big(\bfC \bfG (\bfC \bfG)\tp\big)
  = \tr\big(\bfC\tp \bfC \bfu\bfu\tp)\leq\max\limits_{\bfu\in\{\pm 1\}^k}\ltwo{\bfC\bfu}^2 \le 1. 
\]

\mypar{The rows of $\bfG$ are all orthogonal}
This condition implies $\hbfG = \bfG \bfG\tp$ is a diagonal matrix with diagonal entries in $[0, 1]$, and so \cref{eq:traceCG} implies $\lfrob{\bfC \bfG}^2 \le \tr(\bfX)$. It is thus sufficient to show
\[ \tr(\bfX) \le \max_{\bfu \in \{-1, 1\}^k} \tr(\bfX \bfu \bfu\tp).
\]
We give a construction for a $\bfu$ that shows this. Observe
\begin{align*}
\tr(\bfX \bfu \bfu\tp) 
  &= \tr(\bfX) + 2\sum_{i=1}^k u_i 
     \underbrace{\sum_{j=1}^{i-1} u_j \bfX\idx{i}{j}}_{b_i}. 
\end{align*}
Observe we can choose $u_1 = 1$ and then $u_i = \text{sign}(b_i)$ since $b_i$ depends only on $\bfC$ and the previously fixed $u_j$ for $j < i$, ensuring the double sum on the right is non-negative, completing the proof.



\mypar{$\lfrob{\bfC\bfG} \leq \pi / 2$ without assuming conditions (1)-(4):}

This argument is essentially due to~\citep{denisov_grothendieck}; we inline here for completeness.

We begin by noting that the conditions on $\bfC$ imply that $\bfX := \bfC^\top \bfC$ has $\ell_\infty$-to-$\ell_1$ operator norm 1. This follows from duality. Denote by $\| \cdot \|_{p, q}$ the norm of a matrix as an operator on vectors from $\ell_p$ to $\ell_q$. By assumption,

$$\sup_{\contrib: \|\contrib\|_\infty \leq 1} \|\bfC\contrib\|_2^2 \leq 1,$$
which is to say that $\|\bfC\|_{\infty, 2}  \leq 1$. Duality, therefore, implies that $\|\bfC^\top\|_{2, 1} \leq 1$, and $\|\bfX\|_{\infty, 1} \leq 1$.

Now, for $\bfX = \left(x_{ij}\right),$

\begin{equation}\label{eq:grothendieck_trace_expansion}
\|\bfC\bfG\|_F = \tr\left(\bfX\bfG\bfG^\top\right) = \sum_{i, j=1}^k x_{ij}\langle\bfg_j, \bfg_i\rangle.
\end{equation}

Bounding equation~\cref{eq:grothendieck_trace_expansion} in terms of the $\ell_\infty$-to-$\ell_1$ operator norm of $\bfX$ is the subject of a famous result of Grothendieck~\citep{grothendieck1956resume}; e.g., as stated in~\citep{Lindenstrauss1968}, Grothendieck's inequality may be states precisely as:

\begin{equation}\label{eq:grothendieck_bound}
     \sum_{i, j=1}^k x_{ij}\langle\bfg_j, \bfg_i\rangle \leq K \|\bfX\|_{\infty, 1},
\end{equation}

where $K$ (known as the Grothendieck constant) is independent of $\mdim$, but not known exactly in general.

However, in our case, our matrix $\bfX$ is symmetric by construction; in the case of symmetric $\bfX$ the constant $K$ may be replaced by $\pi / 2$ (see Eq. 43 of \citep{khot2011grothendiecktype}), and this constant is known to be sharp (IE, the best constant which holds for all symmetric matrices $\bfX$ and in all dimensions $\mdim$).

\mypar{Replacing the condition on $\bfG$ to $\forall i \in [k], \|\bfG\idx{i}{:}\|_1 \le 1$, then $\lfrob{\bfC \bfG } \le 1$:} First notice that since $\lfrob{\bfC\bfG}^2$ is a convex function, the maximum happens at the extreme points of the constraint set $\calG = \{\bfG \mid \forall i \in [k], \|\bfG\idx{i}{:}\|_1 = 1\}$. We use \cref{lem:rowstoc} to identify the extreme points of $\calG$.

\begin{claim}[Theorem 1 in~\citet{cao2022centrosymmetric}]
The set of extreme points of the set of $k\times d$ row-stochastic matrices are precisely the set of row permutation matrices, i.e., set of the matrices with entries in $\{0,1\}^{k \times d}$ and each row has exactly one non-zero entry.
\label{lem:rowstoc}
\end{claim}
Notice that the constraint on any matrix $\bfG\in\calG$ is oblivious to the sign, meaning, we can flip the sign of any set of entries in $\bfG$ and the new matrix will still be in $\calG$. This along with \cref{lem:rowstoc} immediately implies that the set $\calH=\{H\in\{-1,0,1\}^{k\times d}: \forall i\in[k], \|\bfH_{[i,:]}\|_0=\linfty{\bfH_{[i,:]}}=1\}$ is the set of extreme points of the set $\calG$. (If the set of extreme points of $\calG$ is larger than $\calH$, then the signs of any such extreme point can be flipped to create a new extreme point of row-stochastic matrices, which would violate \cref{lem:rowstoc}.)

It is not hard to observe that for any $H\in\calH$, there exists an $\bfu_{\bfH}\{\pm 1\}^k$ s.t. $\lfrob{\bfC\bfH}=\ltwo{\bfC\bfu_{\bfH}}$. Since, $\ltwo{\bfC\bfu_{\bfH}}\leq \max\limits_{\bfu\in\{\pm 1\}^k}\ltwo{\bfC\bfu}$ for any choice of $\bfu_{\bfH}$, and the fact that $\max\limits_\bfG\lfrob{\bfC\bfG}$ is reached at one of the matrices in $\calH$, the claim in Theorem~\ref{thm:901r} follows.
\end{proof}

\subsection{A counterexample for general $\bfC$}
\label{sec:counterexample}

Theorem~\ref{thm:901r} indicates the possibility of the following conjecture being true, because of it being true in so many special cases: If $\max\limits_{\bfu\in\{\pm 1\}^k}\ltwo{\bfC\bfu}\leq 1$, then $\lfrob{\bfC\bfG}\leq 1$ for all $\bfG$ with row $\ell_2$-norm at most one. Unfortunately, we show that the conjecture is not true when $k>2$, as shown by the following counterexample with $\dimension = k = 3$ and $\mdim = 2$:%

\[
\bfC = \frac{1}{\sqrt{24}}\begin{bmatrix}
2 & 1 &  1 \\
1 & 2 & -1 \\
1 &-1 &  2 
\end{bmatrix}
\quad \text{and} \quad
\bfG = \frac{1}{\sqrt{5}}\begin{bmatrix}
2 &  1 \\
2 & -1 \\
1 &  2
\end{bmatrix}
\]
Direct calculation shows  $\max_{\bfu\in\{\pm 1\}^k} \ltwo{\bfC\bfu} = 1$, but $\lfrob{\bfC\bfG} = \sqrt{1.1} \approx 1.049$.








\subsection{Proof of \cref{cor:matrix_to_vector_sens}}
\begin{namedcorollary}[\cref{cor:matrix_to_vector_sens}]
When per-step contributions are bounded by $\clipnorm=1$, for any participation pattern $\Pi$ and dimensionality $\mdim \ge 1$, when $\bfC\tp\bfC \ge 0$ elementwise, we have
\[
\sens_{\deltaset^\mdim_\Pi}(\bfC) = \sens_{\deltaset^1_\Pi}(\bfC).
\]
\end{namedcorollary}
\begin{proof}
To see $\sens_{\deltaset^\mdim_\Pi}(\bfC) \ge \sens_{\deltaset^1_\Pi}(\bfC)$, suppose $\bfu^\star = \argmax_{u \in \deltaset^1_\Pi} \norm{\bfC \bfu}_2$, and observe we can construct a $\bfG$ such that $\lfrob{\bfC \bfG} = \norm{\bfC \bfu^\star}_2$ by taking rows $\bfG\idx{i}{:} = (\bfu^\star_i, 0, \dots, 0) \in \R^\mdim$ for $i \in [\dimension]$.

For the other direction, for each $\pi \in \Pi$, we apply \cref{thm:multidim} to the matrix $\bfC_\pi = \bfC\idx{:}{\pi}$, and observe $\bfC\tp\bfC \ge 0$ is sufficient to imply $\bfC\tp_\pi \bfC_\pi \ge 0$.
\end{proof}

\mypar{Note} The condition $\bfC\tp\bfC \ge 0$ is sufficient but not in fact necessary for \cref{cor:matrix_to_vector_sens} to hold. In particular, for \multiepoch-participation $\Pi$, the sub-matrices $\bfC\tp_\pi \bfC_\pi$ for $\pi \in \Pi$ ``touch'' only $\maxpart^2 \assumedsep$ entries of the $\dimension^2 = \maxpart^2 \assumedsep^2$ entries of $\bfC\tp \bfC$; the other entries of $\bfC\tp \bfC$ could in fact be negative. However, we did not need to use this observation for any of the matrices in our experiments.



\subsection{Proof of \cref{thm:eval_sens_ub}}\label{app:ssec:eval_sens_ub_proof}

\begin{namedtheorem}[\cref{thm:eval_sens_ub}]
Let $\bfC \in \R^\dimdim$, and take some participation pattern $\Pi$, with $\maxpart = \max_{\pi \in \Pi} |\pi|$ the maximum number of participations. 
With $\bfC\idx{:}{\pi}$ representing to selecting the columns of the matrix $\bfC$ indexed by $\pi$ and $\ltwo{\cdot}$ the spectral matrix norm, let
$
\lambda=\max\limits_{\pi \in \Pi}\ltwo{\bfC\idx{:}{\pi}}
$. Then 
\[
 \sens_{\deltaset^1_\Pi}(\bfC) \le \lambda\sqrt{k}.
\]
\end{namedtheorem}
\btodo{I think we should be able to extend this to $\sens_{\deltaset^\mdim_\Pi}(\bfC) \le \lambda\sqrt{k}$ using straightforward properties of matrix norms.}

\begin{proof}
By assumption we have $\norm{\bfu} \le \sqrt{k}$, and so
\begin{equation*}
    \max_{\contrib \in \deltaset} \ltwo{\bfC\contrib}
    \leq \max_{\pi \in \Pi}\max_{\contrib \in \deltaset} \ltwo{\bfC[:,\pi]}\ltwo{\contrib}
    \leq \lambda\sqrt{k}.
\end{equation*}
\end{proof}

